\documentclass[12pt, a4paper]{article}

% --- 必備套件 ---
\usepackage[BoldFont, SlantFont]{xeCJK} % 中文字型處理
\usepackage{amsmath, amssymb, amsthm}   % 數學符號與定理環境
\usepackage{geometry}                   % 版面設定
\usepackage{hyperref}                   % 超連結
\usepackage{tikz}                       % 繪圖
\usepackage{tikz-3dplot}                % 3D繪圖輔助

% --- 字型與格式設定 ---
\setCJKmainfont{Noto Sans CJK TC} 
\geometry{margin=1in}
\everymath{\displaystyle} % 避免分數縮小

% --- 定理與練習環境設定 ---
\newtheorem{proposition}{命題}
\newtheorem{exercise}{練習}
\renewcommand{\proofname}{\textbf{證明}}

\begin{document}

\title{《Set》遊戲的數學}
\author{A. Bolotin (А. БОЛОТИН)}
\date{КВАНТ (Kvant), 2025年 第8期}
\maketitle

\section*{前言與遊戲規則}

紙牌遊戲《Set》是由劍橋大學族群遺傳學專家瑪莎·法爾科（Marsha Falco）於 1974 年發明的。她當時正在研究牛的視力缺陷、馬的步態異常以及德國牧羊犬的癲癇問題。在向獸醫解釋遺傳學中出現的組合學問題時，她使用了帶有特殊符號的卡片。後來她意識到，組合這些卡片可以成為一款遊戲的基礎，《Set》就這樣誕生了。

首先，我們將探討《Set》遊戲本身及其相關的基礎組合學問題，然後再過渡到推廣情形與該遊戲相關的未解數學問題。

玩《Set》遊戲需要一副牌。每張牌上印有一、二或三個相同的圖形：菱形、橢圓形或波浪形。牌上的所有圖形均為藍色、紅色或綠色，並具有三種填充類型之一：實心、網底或空心。因此，每張牌都有四個特徵：圖形數量、圖形形狀、顏色和填充，每個特徵可以取三個不同的值。接下來我們將牌的特徵記為以下形式：
$$ \text{(數量, 形狀, 顏色, 填充)} $$

遊戲開始時，桌面上會擺放 12 張牌。玩家的目標是找出能構成「Set」（集合）的三張牌。所謂的「Set」是指這樣的三張牌：在每一個特徵上，它們要麼完全相同，要麼完全不同。例如，牌 (1, 菱形, 藍色, 實心)、(2, 菱形, 紅色, 實心) 與 (3, 菱形, 綠色, 實心) 構成了一個 Set，因為在第一個和第三個特徵（數量和顏色）上它們完全不同，而在第二個和第四個特徵（形狀和填充）上它們完全相同。

當玩家找到一個 Set 時，便將這三張牌收歸己有，並在桌面上補充新的牌。如果玩家們一致認為桌面上已沒有 Set，則會從牌疊中再翻出三張牌放到桌面上。遊戲持續進行，直到牌疊耗盡，且玩家們同意桌面上剩下的牌無法再構成 Set 為止。獲得最多牌的玩家獲勝。

\section*{基礎組合學問題}

第一個自然的問題是：如果牌不重複，這副牌需要多少張？同時，計算總共有多少種三張牌的組合能構成 Set 也是很有趣的課題。以下兩個命題回答了這些問題，但在繼續閱讀之前，建議讀者先自行思考。

\begin{proposition}
這副牌共有 81 張。
\end{proposition}
\begin{proof}
每張牌有 4 個特徵，每個特徵可以取 3 個值，因此總共有 $3^{4}=81$ 張牌。
\end{proof}

\begin{proposition}
總共有 1080 種三張牌的組合可以構成 Set。
\end{proposition}
\begin{proof}
考慮任意兩張牌。請注意，恰好只有一張牌能與這兩張牌構成 Set。考慮任意一個特徵。如果前兩張牌在這個特徵上相同，那麼第三張牌也必須與它們在此特徵上相同。例如，如果前兩張牌是紅色的，那麼第三張牌也必須是紅色的。如果前兩張牌在這個特徵上不同，那麼第三張牌也必須在這個特徵上與前兩者皆不同。例如，如果第一張牌是紅色的，第二張牌是藍色的，那麼第三張牌就必須是綠色的。

因此，任意一對牌都能唯一地擴展成一個 Set，且任何一個 Set 都可以由三對牌生成。相異牌的對數總共有 $C_{81}^{2}$ 對，這意味著 Set 的總數為：
$$\frac{C_{81}^{2}}{3}=\frac{81\cdot80}{3\cdot2}=1080$$
\end{proof}

值得注意的是，三張相異牌的組合總數為：
$$C_{81}^{3}=\frac{81\cdot80\cdot79}{3!}=85320$$
那麼，Set 所佔的比例為 $\frac{1080}{85320}=\frac{1}{79}\approx0.012658$，也就是說，大約每 79 組三張牌中只有一組能構成 Set。

\begin{exercise}
若三張牌在 $a$ 個特徵上完全不同，且在 $b$ 個特徵上完全相同，我們稱該 Set 具有類型 $(a,b)$。例如，組合 (1, 菱形, 藍色, 實心)、(1, 菱形, 紅色, 實心)、(1, 菱形, 綠色, 實心) 具有類型 (1,3)，因為這些牌僅在顏色上不同，而在其他特徵上皆相同。每種可能類型的 Set 共有多少個？每種類型佔所有 Set 的比例各是多少？
\end{exercise}

請注意以下細節：遊戲規則暗示了在桌面上的 12 張牌中，可能找不到任何一個 Set。那麼，需要擺放多少張牌才能保證至少存在一個 Set 呢？我們試圖找出一個足夠大且不包含 Set 的牌集合。考慮一個牌集合，其上的圖形數量為 1 或 2，圖形為菱形或橢圓形，顏色為藍色或紅色，且填充為實心或網底。這樣的牌總共有 $2^{4}=16$ 張，因為每個特徵只能取 2 個值。

可以發現其中沒有三張牌能構成 Set。的確，假設有某三張牌構成了一個 Set。考慮第一個特徵，根據鴿巢原理（抽屜原理），必有兩張牌在該特徵上相同，但如此一來第三張牌也必須與它們在該特徵上相同。對其他特徵進行類似推論，可知這三張牌在所有特徵上都必須相同。但這與所有牌皆相異的前提矛盾。因此，為了確保存在 Set，必須至少擺放 17 張牌。然而，這仍無法確定 17 張牌是否就足夠了。

\begin{exercise}
驗證圖 2 中以黃色標示的卡片構成了一個不包含 Set 的集合。
\end{exercise}

上述練習表明，實際上必須擺放至少 21 張牌。事實證明，這個條件也是充分的：在任意 21 張牌中，總是存在三張能構成 Set 的牌。這個結論是由義大利數學家朱塞佩·佩萊格里諾（Giuseppe Pellegrino）於 1971 年證明的。然而該證明相當困難，因此我們在此省略。請注意，佩萊格里諾甚至在《Set》這款遊戲被發明之前，就已經證明了與其相關的數學結論！這是因為類似的問題早在數學界獨立出現過了。

\section*{數學形式化與幾何詮釋}

現在，讓我們將此問題進行數學形式化。將每個特徵的值以數字 0、1、2 進行編碼。那麼每張牌都可以表示為一個向量 $\vec{x}=(x_{1};x_{2};x_{3};x_{4})$，其中 $x_{i}$ 的取值為 0、1、2，同時 $x_{1}$ 代表數量、$x_{2}$ 代表形狀、$x_{3}$ 代表顏色、$x_{4}$ 代表填充。現在我們可以定義牌的加法。具體來說，我們將代表牌的兩個向量進行逐項模 3 加法：
$$(x_{1};...;x_{4})+(y_{1};...;y_{4})=(x_{1}+y_{1}\pmod{3};...;x_{4}+y_{4}\pmod{3})$$
牌的減法以及牌乘上一個純量也可以類似地定義。

\begin{proposition}
證明卡片 $\vec{x}$、$\vec{y}$、$\vec{z}$ 構成 Set 的充要條件是 $\vec{x}+\vec{y}+\vec{z}=(0;0;0;0)$。請驗證此條件與以下兩個條件分別等價：$\vec{x}+\vec{z}=2\vec{y}$ ； $\vec{y}-\vec{x}=\vec{z}-\vec{y}$。
\end{proposition}
\begin{proof}
假設 $\vec{x}$、$\vec{y}$、$\vec{z}$ 構成 Set。考慮第一個座標 $x_{1}$、$y_{1}$、$z_{1}$。如果它們相同，則 $x_{1}+ y_{1} + z_{1} = 3x_{1} \equiv 0 \pmod{3}$。
如果它們皆不相同，則 $x_{1}+ y_{1} + z_{1} = 0+1+2 \equiv 0 \pmod{3}$。對其他座標可進行類似的推理，即得 $\vec{x}+\vec{y}+\vec{z}=(0;0;0;0)$。

反之，假設 $\vec{x}+\vec{y}+\vec{z}=(0;0;0;0)$。考慮第一個座標。假設其中有相同的座標。不失一般性，我們可以假設 $x_{1}=y_{1}$。那麼 $z_{1}=0-x_{1}-y_{1}= -2x_{1} \equiv x_{1} \pmod{3}$，也就是說向量的所有第一座標都相同。因此，向量的第一座標要麼全部相同，要麼沒有任兩個相同。對其他座標進行類似的推理。而這個條件正是牌構成 Set 的意義。

等價性 $\vec{x}+\vec{y}+\vec{z}=(0;0;0;0)\Leftrightarrow\vec{x}+\vec{z}=2\vec{y}$ 來自於 $-1\equiv2 \pmod{3}$，而等價性 $\vec{x}+\vec{z}=2\vec{y}\Leftrightarrow\vec{y}-\vec{x}=\vec{z}-\vec{y}$ 則是顯而易見的。命題得證。
\end{proof}

請注意，最後一個條件意味著 $\vec{x}$、$\vec{y}$、$\vec{z}$ 構成了一個等差數列。利用剛剛證明的命題，現在可以輕鬆解決以下練習。

\begin{exercise}
在遊戲過程中，玩家們共收集了 26 個 Set。證明剩下的三張牌同樣構成一個 Set。
\end{exercise}

現在想像一下，牌擁有 $n$ 個特徵，而不是 4 個特徵，其中 $n$ 是某個自然數。那麼卡片可以表示為向量 $\vec{x}=(x_{1};x_{2};...;x_{n})$，其中 $x_{i}$ 的取值為 0、1、2。我們為 $n=2$ 的情況提供一個幾何直觀解釋（其他 $n$ 的情況可以類似地討論）。在平面上標出 9 個點，每個點的座標皆取值 0、1、2。因此，每張牌對應於平面上的一個點。命題 3 說明了在這種幾何詮釋中，Set 就等同於「在模 3 意義下」的直線（圖 3）。這種幾何詮釋在後續內容中將會幫助我們。

\begin{center}
\begin{tikzpicture}[scale=1]
    % 圖 3 的重現：3x3 點與模 3 直線
    \foreach \x in {0,1,2}
        \foreach \y in {0,1,2}
            \fill (\x,\y) circle (2.5pt);
    
    % 畫幾條代表 Set 的線
    \draw[thick, blue] (0,0) -- (2,2);
    \draw[thick, red] (0,2) -- (2,0);
    \draw[thick, green!70!black] (0,1) -- (2,1);
    
    \node[below] at (1,-0.5) {圖 3：Set 作為 $n=2$ 時的「直線」};
\end{tikzpicture}
\end{center}

令 $r(n)$ 表示不包含 Set 的最大卡片集合的大小。如上所述，我們已經了解 $r(4)=20$。同樣容易看出 $r(1)=2$、$r(2)=4$（請自行驗證！）。

\begin{exercise}
證明 $r(n)\ge2^{n}$。
\end{exercise}
\begin{exercise}
證明 $r(3)\ge9$。
\end{exercise}

利用幾何詮釋，我們將證明上界 $r(3)\le9$，結合練習 5 即可得到 $r(3)=9$。

\begin{proposition}
對於三個特徵，不包含 Set 的卡片數量不超過 9。
\end{proposition}
\begin{proof}
類似於二維情況，考慮三維空間中座標為 0、1、2 的 27 個點，並假設 $r(3)\ge10$。那麼在這些點中，存在一個由 10 個點組成的集合 $M$，其中沒有任何三點共線（即不構成 Set）。我們注意到整個空間可以被劃分為 3 個平行平面。但由於 $r(2)=4$，每個平面上最多只能包含 4 個屬於 $M$ 的點。因此，$M$ 中的點在這些平面上的分佈只能符合劃分 $10=4+4+2=4+3+3$。

令 $\alpha_{0}$ 為這些平面中包含 2 或 3 個點（不妨設為 $A$、$B$ 以及可能的 $C$）的平面。作一條通過 $A$ 和 $B$ 的直線 $l$。存在恰好三個平面 $\alpha_{1}$、$\alpha_{2}$、$\alpha_{3}$，它們與 $\alpha_{0}$ 交於直線 $l$（圖 4 展示了這些平面以及點 $A$ 和 $B$ 的相對位置範例，請確認一般情況亦同理）。此外，這些平面共同覆蓋了全部的 27 個點。由於 $r(2)=4$，在每個平面 $\alpha_{1}$、$\alpha_{2}$、$\alpha_{3}$ 上，除了 $A$ 和 $B$ 之外，最多只能再有 2 個屬於 $M$ 的點。回想在 $\alpha_{0}$ 上，除了 $A$ 和 $B$ 之外，最多只能再有 1 個屬於 $M$ 的點。因此，$M$ 中的總點數最多為 $2+(1+2+2+2)=9$ 點，產生矛盾。這意味著 $r(3)\le9$。命題得證。
\end{proof}

\begin{center}
\tdplotsetmaincoords{70}{110}
\begin{tikzpicture}[tdplot_main_coords, scale=1.5]
    % 示意圖 4：通過一條直線交會的四個平面
    \draw[thick, red] (0,-1.5,0) -- (0,1.5,0) node[midway, above right, black] {$l$};
    
    \fill[blue, opacity=0.2] (-1.5,-1.5,0) -- (1.5,-1.5,0) -- (1.5,1.5,0) -- (-1.5,1.5,0) -- cycle;
    \fill[green, opacity=0.2] (0,-1.5,-1.5) -- (0,-1.5,1.5) -- (0,1.5,1.5) -- (0,1.5,-1.5) -- cycle;
    \fill[yellow, opacity=0.2] (-1.5,-1.5,-1.5) -- (1.5,-1.5,1.5) -- (1.5,1.5,1.5) -- (-1.5,1.5,-1.5) -- cycle;
    \fill[red, opacity=0.2] (-1.5,-1.5,1.5) -- (1.5,-1.5,-1.5) -- (1.5,1.5,-1.5) -- (-1.5,1.5,1.5) -- cycle;
    
    \node[below] at (0,-2,0) {圖 4：通過直線交會的各個平面示意};
\end{tikzpicture}
\end{center}

\section*{加法組合學與未解問題}

因此，我們已經了解到 $r(1)=2$、$r(2)=4$、$r(3)=9$、$r(4)=20$。目前也已知 $r(5)=45$、$r(6)=112$。令人驚訝的是，即使是 $r(7)$ 的值，至今仍未被計算出來！這可能是因為人們主要對 $r(n)$ 在 $n$ 較大時的漸近行為感興趣。而由於精確計算 $r(n)$ 相當困難，數學家轉而尋求建立下界和上界。

基於我們剛剛建立的等式 $r(3)=9$，可以得到一個針對 $r(n)$ 的估計值，當 $n$ 足夠大時，該估計優於練習 4 中給出的結果。令 $k=[\frac{n}{3}]$，即 $\frac{n}{3}$ 的整數部分。令 $A=\{a_{1},...,a_{9}\}$ 為一個由 9 張具有 3 個特徵的牌所組成的集合（如練習 5 中所建構），其中不包含 Set。現在考慮以下具有 $n$ 個特徵的牌集合：將所有特徵每 3 個分作一個區塊，並在每個區塊中放入某個 $a_{i}$ 的特徵值。將最後的 $n-3k$ 個特徵設為 0。例如，當 $n=8$ 時，所建構的集合由形式為 $(\vec{a}_{i},\vec{a}_{j},0,0)$ 的卡片組成。容易看出，在我們建構的集合中沒有三張牌構成 Set，並且該集合總共包含 $9^{k}$ 張牌。因此，
$$r(n)\ge9^{k}=9^{[n/3]}\ge9^{n/3-1}=\frac{1}{9}(\sqrt[3]{9})^{n}\approx\frac{1}{9}\cdot2.08^{n}$$

不難驗證，當 $n\ge56$ 時，由此得到的估計強於練習 4 中的結果。也就是說，對於較大的 $n$，只有指數項 $(\sqrt[3]{9})^{n}$ 是重要的，常數 $\frac{1}{9}$ 並不具有決定性作用。這也是為何在推導估計界限時，通常只關注底數。利用 $r(4)$、$r(5)$、$r(6)$ 的值以及類似的推理，可以進一步改進這個下界。

\begin{exercise}
基於 $r(4)=20$，推導出 $r(n)$ 的下界達到 $(\sqrt[4]{20})^{n}\approx2.11^{n}$ 的量級。
\end{exercise}

目前已知的最佳下界約為 $2.2202^{n}$ 的量級。除了下界，上界同樣引人關注。正如我們所見，下界的形式為 $c^{n}$，其中 $c$ 是一個常數。當然，顯然的上界為 $r(n)\le3^{n}$（因為共有 $n$ 個特徵，每個特徵取 3 個值），也具有相同的形式。然而在 1987 年，有人提出了一個猜想：$r(n)\le c^{n}$，其中 $c<3$。傑出的數學家陶哲軒（Terence Tao）曾稱這個猜想是他最喜愛的未解問題。早期得到的其中一個上界是：
$$r(n)\le\frac{2\cdot3^{n}}{n}$$

請讀者自行驗證，這樣的估計比形如 $c^{n}$（$c<3$）的界限還要弱。這個界限曾數次被改進，直到 2016 年，上述猜想才被徹底證明。目前已知的最佳上界約為 $2.756^{n}$ 的量級。關於 $r(n)$ 準確漸近行為的問題，至今仍然是開放的。

我們所探討的這個問題是加法組合學（Additive Combinatorics）中的經典課題。直觀而言，加法組合學致力於尋找組合結構與加法結構之間的聯繫，有時也會將乘法結構納入其中。這些概念並非絕對嚴謹，但為了理解其內涵，我們來看一個例子。自然數集合是「加法性」定義的：其第一個元素是 1，之後的每個元素都是前一個元素加上 1 得到的。另一方面，根據算術基本定理，每一個大於 1 的自然數都可以唯一分解成質數的乘積。因此，在某種意義上，質數構成了自然數的乘法結構，而每個數的大小（加 1 的次數）則構成了加法結構。當混合加法和乘法結構時，經常會產生複雜的問題，數學家們為此投入了多年的心血。例如，會出現以下問題：

\textbf{關於乘法結構的加法問題：}
\begin{itemize}
    \item 第 $n$ 個質數 $p_{n}$ 的大小如何變化？
    \item 關於相鄰質數之間的差值 $p_{n+1}-p_{n}$，我們能得出什麼結論？
\end{itemize}
第一個問題與質數分佈定理以及黎曼猜想（少數未解決的希爾伯特問題之一）有關。關於第二個問題，直到 2013 年才證明了存在一個常數 $C$，使得對於無限多對相鄰質數滿足 $p_{n+1}-p_{n}\le C$。後來證明 $C$ 可以取為 246。著名的孿生質數猜想則認為 $C=2$ 也是成立的。

\textbf{關於加法結構的乘法問題：}
\begin{itemize}
    \item 如何將數字 $n$ 分解為質因數？
    \item 數字 $n$ 與 $n+1$ 的質因數分解之間有何關聯？
\end{itemize}
第一個問題涉及演算法理論。目前尚不知道任何「足夠快」的質因數分解演算法。不過，現已有 AKS 質數測試（Agrawal-Kayal-Saxena）能夠「足夠快」地檢驗一個數是否為質數，以及著名的秀爾演算法（Shor's algorithm），這是一種能在量子電腦上執行的「足夠快」質因數分解演算法。第二個問題則與考拉茲猜想（Collatz conjecture）有關，保羅·艾狄胥（Paul Erdős）曾對此評論道：數學界還沒有準備好應對這類問題。

《Set》遊戲最初是作為一個組合對象出現的，我們討論的是具有取不同數值的若干特徵的卡片。然而，命題 3 表明，關於 Set 構成的問題完全可以用加法結構的術語來提出。另一方面，我們也看見了一個純粹的幾何詮釋：Set 就是某個特定空間中的直線。在一個簡單的兒童遊戲背後，竟然隱藏著如此豐富且多元的數學世界！

\end{document}